\section{Introduction}

The dorsal hippocampus (dHPC) occupies a central position in the vertebrate cognitive map, transforming sensory and self-motion information into spatial representations that guide flexible behavior \citep{OKeefe1971PlaceCells,Skaggs1993Information}. Decades of electrophysiology and, more recently, large-scale two-photon calcium imaging have shown that CA1 pyramidal neurons encode not only location but also velocity, reward, and contextual variables relevant to ongoing task demands \citep{Gauthier2018Reward,Sosa2018Reward,Sato2020Overrepresentation}. How this rich representational repertoire is transformed into adaptive action, however, depends on where in the brain hippocampal activity is read out. A key downstream target is the nucleus accumbens (NAc), long proposed to serve as a limbic-motor interface that translates motivationally salient signals into goal-directed behavior \citep{Mogenson1980Motivation,Britt2012Synaptic}.

A central question is whether the dHPC projection to NAc carries a stereotyped copy of the hippocampal code or, instead, a specialized subset of signals tailored to reward-guided action. Anatomical and functional evidence suggests that hippocampal projection neurons are not uniform: ventral CA1 neurons projecting to different downstream regions support distinct behaviors such as anxiety, fear, and reward \citep{Ciocchi2015Selective,Bakoyiannis2023Pathway}. Recent work has further shown that the dHPC-NAc pathway is capable of driving reinforcement, engaging accumbal dynorphin circuits and integrating with prefrontal inputs to modulate engagement in reward-seeking \citep{Ibrahim2024DorsalHPC,Iyer2025Engagement}. Within the NAc itself, dopamine-dependent mechanisms help organize memory modules that support goal-directed navigation \citep{Jung2024Dopamine}. Yet despite this progress, the population-level content and projection specificity of dHPC information arriving at NAc during actual goal-directed behavior remains largely uncharacterized.

Bridging this gap requires simultaneous access to large populations of hippocampal neurons and a molecularly defined subset that unambiguously projects to NAc, together with behaviorally grounded manipulations of the same pathway. Dual-color two-photon imaging, combined with retrograde viral labeling, now makes such dissection feasible in head-fixed mice performing structured spatial reward tasks \citep{Tervo2016AAVretro,Chen2013GCaMP6,Dana2014Thy1GCaMP6}. Earlier observations of place field overrepresentation near rewarded locations and of conjunctive reward-place coding in distal CA1 \citep{Hok2007Goal,Mamad2017PlaceField,Xiao2020Conjunctive,Kaufman2020LocusCoeruleus} hint that dHPC-NAc neurons may be enriched for such task-aligned representations, but direct, projection-specific evidence has been missing.

Here, we test this hypothesis by combining projection-targeted imaging with optogenetic circuit manipulation during goal-directed appetitive behavior. Using dual-color two-photon microscopy in mice co-expressing pan-neuronal GCaMP6s and retrogradely-driven mCherry in NAc-projecting dHPC neurons, we monitored 5,372 neurons across 19 sessions from 6 mice as they learned a hidden-reward treadmill task. We compared projection-defined (dHPC$\rightarrow$NAc) and non-labeled (dHPC$-$) populations with respect to place coding, reward-zone overrepresentation, speed modulation, appetitive licking, and conjunctive coding. We then asked whether activation of the same projection is causally sufficient to elicit the consummatory behaviors it represents. Our results, summarized below, describe the dHPC-NAc channel as a specialized carrier of reward-anchored spatial and behavioral-state information that is causally engaged in appetitive action.

\section{Related Work}

\paragraph{Hippocampal representations of reward location.}
Reward has long been known to shape the hippocampal map. Place fields tend to cluster near rewarded locations in both freely-moving rodents and head-fixed paradigms, a phenomenon proposed to reflect experience-dependent biasing of the cognitive map by motivationally relevant outcomes \citep{Hok2007Goal,Mamad2017PlaceField,Sato2020Overrepresentation}. Two recent imaging studies identified dedicated populations of reward-responsive cells in CA1 that are distinct from classical place cells and whose activity tracks reward identity and value \citep{Gauthier2018Reward,Sosa2018Reward}. However, the degree to which reward-related clustering and conjunctive reward-place coding generalize to projection-defined subpopulations has remained unclear, and not all tasks produce place-cell sensitivity to goal value \citep{Duvelle2019PlaceCells}. Neuromodulatory systems, including locus coeruleus noradrenergic input, have been implicated in the reorganization of CA1 place cells during reward learning \citep{Kaufman2020LocusCoeruleus}, and conjunctive reward-place coding has been quantified in distal CA1 \citep{Xiao2020Conjunctive}.

\paragraph{The hippocampus-accumbens pathway.}
Anatomical and electrophysiological studies have long identified hippocampal projections to NAc as a key substrate for reward-guided behavior and memory-motivation interactions \citep{Mogenson1980Motivation,Britt2012Synaptic}. Ventral CA1 projections have been shown to route distinct information streams to different downstream targets \citep{Ciocchi2015Selective}, and projection-specific manipulations reveal diverse behavioral roles for these ventral pathways \citep{Bakoyiannis2023Pathway}. More recent work has shifted attention toward the \emph{dorsal} hippocampus-NAc route. A hippocampus-accumbens tripartite motif was shown to guide appetitive memory in space \citep{Trouche2019Place}, and dHPC-NAc terminals are capable of driving reinforcement through accumbal dynorphin neurons \citep{Ibrahim2024DorsalHPC}. Prefrontal and ventral-hippocampal inputs to NAc cooperatively modulate engagement \citep{Iyer2025Engagement}; cholecystokinin interneurons in ventral hippocampus gate contextual reward memory \citep{Nguyen2024Cholecystokinin}; and NAc itself hosts dopamine-dependent memory modules for navigation \citep{Jung2024Dopamine}. Still missing from this picture is a population-scale description of what dHPC-NAc neurons encode during ongoing goal-directed behavior, together with a causal test of whether the same projection drives the behaviors it represents.

\paragraph{Technical advances enabling projection-specific imaging.}
Progress in projection-specific population imaging has been enabled by three converging tools. High-efficiency retrograde AAV variants (e.g., AAVrg) allow Cre delivery to axon terminals with minimal toxicity \citep{Tervo2016AAVretro}. Ultrasensitive calcium indicators such as GCaMP6s, in both AAV and transgenic (Thy1) forms, provide stable signals in large populations over extended recording sessions \citep{Chen2013GCaMP6,Dana2014Thy1GCaMP6}. Finally, robust analysis pipelines for motion correction and source extraction from two-photon data---NoRMCorre and CaImAn---enable unbiased, scalable processing of chronic imaging data sets \citep{Pnevmatikakis2017NoRMCorre,Giovannucci2019CaImAn}. Our study builds directly on these foundations to interrogate, for the first time at this scale, projection-defined dHPC-NAc coding during a structured goal-directed task.
